% hori.tex - 横排 clreq 管道（H1–H4）回归测试
% 覆盖：CJK 断点注入、行首行尾禁则（penalty）、中西间距 1/4em 及其例外、
%       符号分离禁则（数字串/数字+单位/货币）、两字宽标点不可拆、
%       H2 行内优先级二次分配与行末标点半字宽、
%       H3 行间标号（专名号/书名号甲式/着重号）、H4 行间注（拼音/对照）
\documentclass[12pt]{article}
\usepackage[paperwidth=10cm, paperheight=16cm, margin=1.2cm]{geometry}
% 中国大陆偏靠式标点（style=mainland 默认）需要 SC 字面：思源宋体标点靠始端、
% 空白在末端；TW-Kai 的标点为台式居中字面，不能用于中国大陆式测试。
% Script=CJK + Language 必须显式指定：否则 luaotfload 按 latn/dflt 起排，
% 该语言系统的 locl 会把省略号替换为贴基线的西文形（cid63070）；
% hani/ZHS 下省略号保持居中默认形，破折号经 locl 升级为满格中文形。
\usepackage{fontspec}
\setmainfont{Source Han Serif SC}[Script=CJK, Language={Chinese Simplified}]
\usepackage{luatex-cn-hori}
\pagestyle{empty}

% 用例 ID：往内容流写一段 marked content 把用例包起来，让 clreq_test.py
% 能按 ID 定位，而不是靠「页面上第一个含某字的行」——文字选择器一旦不唯一
% （说明文字里也出现同样的字），断言会静默量错对象。
% 安全性：literal 是零宽 whatsit，不构成断点；横排管道把 whatsit 当透明节点
% （hori-pipeline.lua 的 TRANSPARENT），前后相邻字形的判定不受影响，版面不变。
% 用法限制：只包整段/整个用例。塞进正被量间隙的两个字形之间会切断西文
% kern 串（¥1280、Hello 这类），反而扰动被测对象。
% \leavevmode 不可省：\pdfextension literal 不开启水平模式，在竖直模式下
% 它会作为 whatsit 落进主竖直表、留在上一页页底，正文另起段落流到下一页——
% BDC 与它标记的字形就此被分页拆散，断言只会报「PDF 里没有这个用例」。
\newcommand\用例[2]{%
  \leavevmode
  \pdfextension literal page{/Span <</T (#1)>> BDC}%
  #2%
  \pdfextension literal page{EMC}%
}

\begin{document}

% 版面上渲染配置说明：让看输出的人不必翻源码即可知道启用了什么
{\small 【测试配置】本页为宏包默认选项：中国大陆式标点、基本级禁则、行末标点半字宽挤压；未启用引号体例转换与行尾点号悬挂。}

% --- 基本断行与行首行尾禁则：长句在窄版面上多次换行，
%     点号不应出现在行首、开引号不应出现在行尾
天地玄黄，宇宙洪荒。日月盈昃，辰宿列张。寒来暑往，秋收冬藏。闰余成岁，律吕调阳。云腾致雨，露结为霜。金生丽水，玉出昆冈。剑号巨阙，珠称夜光。果珍李柰，菜重芥姜。

子曰：「学而时习之，不亦说乎？有朋自远方来，不亦乐乎？人不知而不愠，不亦君子乎？」又曰：「温故而知新，可以为师矣。」

% --- 中西混排：汉字与西文、数字之间应有约 1/4 字宽的间距
本项目基于 LuaTeX 引擎开发，版本号为 1.18.0，托管在 GitHub 平台上。使用 expl3 编程层与 Lua 脚本协同工作，支持 OpenType 字体特性。

% --- 中西间距例外：点号与夹注号旁不加间距
他说：Hello 世界！括号（English）内外、标点，abc 之后。

% --- 符号分离禁则：数字串、数字+单位、货币不拆行
\用例{sep-spaced}{测量结果为 1234 毫米，占比 95\%，温度 37℃，误差 ±3，价格 ¥1280，共计 100 件样品，编号从 001 到 999，全部合格。}

% --- 同上但源码不加空格：中西间距 1/4em 应由管道自动插入
\用例{sep-unspaced}{复测结果为5678毫米，占比87\%，温度36℃，误差±2，价格¥999，共计50件样品，标号从010到099，使用LuaTeX编译，全部达标。}

% --- 中西间距的「中文一侧」到底包括谁。两条各自单独成段：单行段落的末行
%     不参与两端对齐，量到的就是管道插入的原值，不会被 H2 的行内挤压改写
%     （整段排满时 cjk_western 是第一优先级的可挤对象，会被压没）。
% - / · – 的字面是西文自带的半角标点，夹在西文里不构成中西边界，不加：
\用例{west-form-punct}{1/4 em、a·b、TW-Kai。}

% 破折号、省略号只有中文用、字面满幅，是中文的一部分，照加：
\用例{cjk-punct-western}{破折号——abc 省略号……xyz。}

% --- 两字宽标点：破折号与省略号
生活就像一盒巧克力——你永远不知道下一颗是什么味道……这就是命运给出的安排。

% --- 连续标点缩减与段首缩进行首开括号：段落以「开头（始侧空白应被缩减，
%     缩进视觉上仍为两字），段内含 」。/？」/：「 等连排组合
「大学之道，在明明德，在亲民，在止于至善。」此句出自《大学》（儒家经典）；所谓「知止而后有定」，即此理也。

% --- H3/H4 段落带行间标注（单面装），按 clreq 行距须 ≥ 0.5em 行间空，
%     否则包会给出 warning
\begingroup
\linespread{1.6}\selectfont

{\small 【以下两段】行距放大 1.6 倍（行间标注要求行距不小于半字），演示专名号、书名号甲式、着重号，以及拼音、对照行间注。}

% --- H3 行间标号：专名号（直线）、书名号甲式（波浪线）、着重号（点）
%     相邻专名各自分开；着重号不加在标点上
太史公\专名{司马迁}著\书名号甲{史记}，首篇为\书名号甲{五帝本纪}。其文\着重{究天人之际，通古今之变}，向为史家推重。\专名{汉}\专名{高祖}起于\专名{沛县}，事见本纪。

% --- H4 行间注：拼音标音（分词连写）与中外文对照
\用例{h4-ruby}{\拼音{中国}{zhōng guó}的\拼音{汉字}{hàn zì}历史十分悠久，排版时可以使用\拼音{拼音}{pīn yīn}为生僻的文字标注读音；\对照{印刷}{printing}术的发明，对于文化的传播意义重大。}

\endgroup

\end{document}
